% Part: first-order-logic
% Chapter: semantics
% Section: intro-semantics

\documentclass[../../../include/open-logic-section]{subfiles}

\begin{document}

\olfileid{fol}{syn}{its}

\olsection{పరిచయం}

వ్యక్తీకరణలకు అర్థం ఇవ్వడం అర్థవిచార పరిధి. అర్థవిచారంలో కేంద్ర
భావన, ఒక \tetoken{నిర్మాణంలో}{!!a{structure}} సంతృప్తి. \tetoken{ఒక నిర్మాణం}{!!^a{structure}} భాషలోని
మౌలిక భాగాలకు అర్థం ఇస్తుంది: \tetoken{ఒక వ్యక్తి క్షేత్రం}{!!a{domain}} అనేది ఖాళీ కాని వస్తువుల
సమితి. పరిమాణకాలు ఈ వ్యక్తి క్షేత్రమంతటా పరిధి కలిగి ఉన్నట్లు
అర్థం చేసుకుంటాం; \tetoken{స్థిరసంకేతాలకు}{!!{constant}s} వ్యక్తి క్షేత్రంలోని మూలకాలను,
\tetoken{ప్రమేయ సంకేతాలకు}{!!{function}s} \tetoken{వ్యక్తి క్షేత్రం}{!!{domain}} నుంచి అదే వ్యక్తి క్షేత్రానికి ప్రమేయాలను,
\tetoken{విధేయ సంకేతాలకు}{!!{predicate}s} \tetoken{వ్యక్తి క్షేత్రంపై}{!!{domain}} సంబంధాలను నిర్దేశిస్తాం. మౌలిక
పదజాలానికి చేసిన ఈ నిర్దేశాలతో కలిసి \tetoken{వ్యక్తి క్షేత్రం}{!!{domain}} ఒక \tetoken{నిర్మాణాన్ని}{!!a{structure}}
ఏర్పరుస్తుంది. \tetoken{చరాలు}{!!^{variable}s}, \tetoken{సూత్రాలలో}{!!{formula}s} కనిపించవచ్చు; కాబట్టి
అర్థవిచారం ఇవ్వడానికి వాటికి కూడా \tetoken{మూలకాలను}{!!{element}s} \tetoken{వ్యక్తి క్షేత్రం}{!!{domain}} నుంచి
నిర్దేశించాలి---ఇదే చర నిర్దేశం. చివరగా, సంతృప్తి సంబంధం వీటన్నిటినీ
కలుపుతుంది. ఒక \tetoken{సూత్రం}{!!^a{formula}}, \tetoken{ఒక నిర్మాణం}{!!a{structure}}~$\Struct{M}$లో
ఒక \tetoken{చర}{!!a{variable}} నిర్దేశం~$s$కు సాపేక్షంగా సంతృప్తమవవచ్చు; దీనిని
$\Sat{M}{!A}[s]$గా రాస్తాం. ఈ సంబంధాన్ని కూడా~$!A$ నిర్మాణంపై
ఆగమనం ద్వారా నిర్వచిస్తాం. ఉదాహరణకు, తార్కిక సంయోజకాల సత్య
పట్టికలను ఉపయోగించి $(!A \land !B)$ సంతృప్తిని $!A$, ~$!B$ల
సంతృప్తి (లేదా అసంతృప్తి) ద్వారా నిర్వచిస్తాం. \tetoken{చర}{!!{variable}} నిర్దేశం
అసంబద్ధమని తరువాత తెలుస్తుంది, ఒకవేళ \tetoken{సూత్రం}{!!{formula}}~$!A$ ఒక
\tetoken{వాక్యం}{!!a{sentence}} అయితే, అంటే దానిలో స్వేచ్ఛా చరాలు లేకపోతే. అందువల్ల \tetoken{వాక్యాలు}{!!{sentence}s},
\tetoken{నిర్మాణాలలో}{!!{structure}s} సంతృప్తమవుతాయని (లేదా కావని) నేరుగా చెప్పవచ్చు.

\tetoken{వాక్యాలు}{!!{sentence}s} కోసం సంతృప్తి సంబంధం~$\Sat{M}{!A}$ ఆధారంగా
చెల్లుబాటుతనం, అర్థపర అనుగమనం, సంతృప్తిపరచదగినతనం అనే మౌలిక
అర్థపర భావనలను నిర్వచించవచ్చు. ఒక \tetoken{వాక్యం}{!!^a{sentence}}, ప్రతి
\tetoken{నిర్మాణం}{!!{structure}} దానిని సంతృప్తిపరిస్తే చెల్లుబాటు అవుతుంది,
$\Entails !A$. ఒక \tetoken{వాక్యాల}{!!{sentence}s} సమితి నుంచి అది అర్థపరంగా
అనుగమిస్తుంది, $\Gamma \Entails !A$, ఒకవేళ ఆ సమితి~$\Gamma$ను
సంతృప్తిపరిచే ప్రతి \tetoken{నిర్మాణం}{!!{structure}}, అందులోని అన్ని \tetoken{వాక్యాలతో}{!!{sentence}s}
పాటు ~$!A$ను కూడా సంతృప్తిపరిస్తే. ఒక \tetoken{వాక్యాల}{!!{sentence}s} సమితిని, ఏదో ఒక
\tetoken{నిర్మాణం}{!!{structure}} అందులోని అన్ని \tetoken{వాక్యాలను}{!!{sentence}s} ఒకేసారి సంతృప్తిపరిస్తే ఆ సమితి
సంతృప్తిపరచదగినది. \tetoken{సూత్రాలు}{!!{formula}s} ఆగమనాత్మకంగా నిర్వచించబడ్డాయి;
\tetoken{సూత్రాలు}{!!{formula}s} నిర్మాణంపై ఆగమనం ద్వారా సంతృప్తి కూడా నిర్వచించబడింది. కాబట్టి
మన అర్థవిచార ధర్మాలను నిరూపించడానికి, నిర్వచించిన అర్థపర భావనలను
పరస్పరం అనుసంధానించడానికి ఆగమనాన్ని ఉపయోగించవచ్చు.

\end{document}
